\chapter{Fazit}

Im Rahmen dieser Arbeit wurde ein Flachlautsprecher auf Basis der DML-Technologie entwickelt, ausgelegt, simulativ untersucht, konstruktiv umgesetzt und messtechnisch bewertet. Ziel war es, die grundsätzliche Eignung eines durch Körperschallwandler angeregten Platte als Lautsprechersystem zu untersuchen und dabei die wesentlichen Einflussgrößen auf das akustische Verhalten herauszuarbeiten. Die Arbeit zeigt, dass DML-Lautsprecher aufgrund ihrer flachen Bauform, der einfachen konstruktiven Integration und der großflächigen Schallabstrahlung ein interessantes alternatives Lautsprecherkonzept darstellen. Insbesondere dort, wo klassische Lautsprecher aufgrund ihrer Bautiefe, ihres Gehäusevolumens oder ihrer gerichteten Abstrahlung nur eingeschränkt eingesetzt werden können, bieten DML-Systeme deutliche konstruktive Vorteile.

Die Untersuchung hat jedoch gezeigt, dass die Entwicklung eines leistungsfähigen DML-Lautsprechers deutlich komplexer ist als der einfache mechanische Aufbau zunächst vermuten lässt. Das akustische Verhalten wird nicht allein durch den verwendeten Exciter bestimmt, sondern maßgeblich durch das Zusammenspiel aus Panelmaterial, Plattengröße, Biegesteifigkeit, Flächenmasse, Randlagerung und Exciter-Position. Da der Schall bei einem DML-System durch die Anregung von Biegewellen und Eigenmoden entsteht, hängt die Qualität der Abstrahlung unmittelbar davon ab, welche Moden im relevanten Frequenzbereich angeregt werden. Eine ungünstige Positionierung des Exciters kann dazu führen, dass einzelne Moden nur schwach oder gar nicht angeregt werden, während andere Moden überbetont werden. Dies führt zu Einbrüchen, Resonanzspitzen und einer insgesamt ungleichmäßigen Wiedergabe.

Daraus ergibt sich, dass eine simulationsgestützte Untersuchung der Modenstruktur für die Auslegung eines DML-Lautsprechers von zentraler Bedeutung ist. Die Simulation ermöglicht es, Eigenformen und Schwingungsbereiche der Platte sichtbar zu machen und daraus geeignete Anregerpositionen abzuleiten. Damit stellt sie einen wichtigen Entwicklungsschritt dar, um eine möglichst effiziente und gleichmäßige Abstrahlung zu erreichen. Eine rein empirische Positionierung des Exciters wäre zwar grundsätzlich möglich, würde jedoch einen erheblich höheren experimentellen Aufwand verursachen und könnte nur schwer auf andere Panelgeometrien oder Materialien übertragen werden. Die Ergebnisse dieser Arbeit verdeutlichen daher, dass die Simulation der Moden keine ergänzende Maßnahme, sondern eine Voraussetzung für die gezielte Entwicklung eines DML-Systems ist.

Kritisch zu bewerten ist außerdem, dass ein DML-Lautsprecher aufgrund seines physikalischen Wirkprinzips keinen von sich aus linearen Frequenzgang bereitstellt. Die Schallabstrahlung entsteht aus der Überlagerung vieler modaler Schwingungen, deren Amplituden und Phasenlagen frequenzabhängig variieren. Dadurch entstehen im Frequenzgang typische Unregelmäßigkeiten in Form von Pegelüberhöhungen und Einbrüchen. Besonders bei einem Fullrange-System ohne Frequenzweiche wird diese Problematik deutlich, da eine einzige Platte den gesamten hörbaren Frequenzbereich wiedergeben soll. Während dies konstruktiv einfach ist, stellt es akustisch einen Kompromiss dar. Vor allem der Tieftonbereich kann aufgrund der begrenzten Luftverdrängung und der geringen Hubfähigkeit der Plattenfläche nur eingeschränkt wiedergegeben werden kann.

Aus diesem Grund ist bei einem DML-Fullrange-System eine nachträgliche Equalisation des Signals notwendig. Durch eine gezielte Entzerrung können Resonanzüberhöhungen reduziert, Pegelunterschiede teilweise ausgeglichen und der subjektive Höreindruck verbessert werden. Die Equalisation kann jedoch die physikalischen Grenzen des Systems nicht vollständig aufheben. Insbesondere eine mangelnde Tieftonabstrahlung kann nur begrenzt elektronisch kompensiert werden, da eine starke Anhebung tiefer Frequenzen zu höherer mechanischer und elektrischer Belastung des Exciters führt. Die Equalisation ist daher als notwendige Optimierungsmaßnahme zu verstehen, ersetzt jedoch keine geeignete mechanische und akustische Grundauslegung.

Zusammenfassend lässt sich feststellen, dass das entwickelte DML-System die grundsätzliche Funktionsfähigkeit der Technologie erfolgreich nachweist. Der Aufbau zeigt, dass mit vergleichsweise einfachen Mitteln ein flacher Lautsprecher realisiert werden kann, der insbesondere durch seine breite Abstrahlung und seine Integrationsfähigkeit überzeugt. Gleichzeitig zeigt die Arbeit aber auch die Grenzen eines einfachen Fullrange-DML-Systems auf. Für eine hohe Klangqualität sind eine gezielte Materialauswahl, eine geeignete Plattengeometrie, eine simulationsgestützte Exciter-Positionierung sowie eine elektronische Entzerrung erforderlich. DML-Lautsprecher eignen sich daher besonders für Anwendungen, bei denen flache Bauweise, Designintegration und großflächige Beschallung im Vordergrund stehen. Für Anwendungen mit hohen Anforderungen an Linearität, Tieftonwiedergabe, maximale Pegelfestigkeit und präzise Lokalisation bleibt hingegen weiterer Optimierungsbedarf bestehen.Damit bildet die Arbeit eine solide Grundlage für die weitere Entwicklung und Optimierung flacher, integrierbarer Lautsprechersysteme auf Basis der DML-Technologie.


\chapter{Ausblick}

Aufbauend auf den Ergebnissen dieser Arbeit ergeben sich verschiedene Ansätze zur Weiterentwicklung des entwickelten DML-Flachlautsprechers. Da ein reines Fullrange-DML-System aufgrund seiner physikalischen Eigenschaften insbesondere im Tieftonbereich und hinsichtlich eines gleichmäßigen Frequenzgangs Einschränkungen besitzt, erscheint eine Erweiterung des Systems sinnvoll.
Ein möglicher nächster Schritt ist die Einbindung eines Subwoofers. Dieser könnte den tieffrequenten Bereich übernehmen, in dem DML-Panels aufgrund ihrer geringen Luftverdrängung nur eingeschränkt effizient arbeiten. Das DML-Panel könnte dadurch gezielt für den Mittel- und Hochtonbereich eingesetzt werden. Durch eine geeignete Trennfrequenz ließe sich die Klangqualität verbessern und die Belastung des Panels im Tieftonbereich reduzieren.
Darüber hinaus bietet sich der Umbau zu einem DSP-gesteuerten Mehrwegesystem an. Ein digitaler Signalprozessor könnte das Audiosignal in unterschiedliche Frequenzbereiche aufteilen, Laufzeitunterschiede korrigieren und den Frequenzgang gezielt entzerren. Besonders ein 3-Wege-System aus Subwoofer, DML-Panel und separatem Hochtonsystem könnte die jeweiligen Stärken der einzelnen Komponenten besser nutzen. Eine solche Weiterentwicklung könnte im Rahmen einer Studienarbeit oder Bachelorarbeit systematisch untersucht, aufgebaut und messtechnisch bewertet werden.
Ein weiterer Ansatzpunkt liegt in der Verbesserung des Simulationsmodells. Zukünftige Simulationen sollten die Sandwichplatte nicht nur als monolithische Ersatzplatte betrachten, sondern den realen Schichtaufbau genauer abbilden. Zusätzlich könnte ein Exciter-Ersatzmodell integriert werden, das Masse, Steifigkeit, Dämpfung, Kontaktfläche und elektromechanische Eigenschaften des Körperschallwandlers berücksichtigt. Dadurch ließe sich das Zusammenspiel zwischen Exciter und Panel realitätsnäher untersuchen. Da ein solches Modell einen deutlich höheren Rechenaufwand verursacht, sollte die Berechnung auf einem leistungsstarken PC oder einer geeigneten Workstation durchgeführt werden. Auf diese Weise könnten Exciter-Positionen, Plattengrößen, Materialien und Lagerungsarten bereits vor dem Bau weiterer Prototypen genauer bewertet werden.

Neben der technischen Optimierung bietet auch die praktische Anwendung an der Hochschule Potenzial. Aufgrund ihrer flachen Bauform und breiten Abstrahlung könnten DML-Lautsprecher als Decken- oder Wandpaneele in Hörsälen, Seminarräumen oder Laboren eingesetzt werden. Besonders als Deckenpanel könnten sie eine gleichmäßige Beschallung ermöglichen und sich unauffällig in die Raumstruktur integrieren. Gleichzeitig könnten die Paneele gestalterische oder informative Funktionen übernehmen.
Für zukünftige Untersuchungen wäre daher zu prüfen, wie sich DML-Paneele in realen Hochschulräumen akustisch verhalten. Dabei sollten Frequenzgang, Sprachverständlichkeit, räumliche Gleichmäßigkeit und Raumakustik betrachtet werden. Insgesamt bietet diese Arbeit damit eine Grundlage für weitere studentische Projekte, in denen Subwoofer-Integration, DSP-Entzerrung, Mehrwegekonzepte, verbesserte Simulationen und praktische Anwendungen von DML-Systemen weiter untersucht werden können.